\begin{tikzpicture}[
  font=\scriptsize,
  execute at begin node={\hyphenpenalty=10000\exhyphenpenalty=10000},
  tanya/.style={draw, rounded corners, fill=orange!12, text width=2.45cm, minimum height=0.95cm, align=center, font=\tiny},
  kelompok/.style={draw, rounded corners, fill=praditagreen!15, text width=2.45cm, minimum height=0.55cm, align=center, font=\scriptsize\bfseries},
  halaman/.style={draw, rounded corners, fill=blue!8, text width=2.45cm, minimum height=1.6cm, align=left, font=\tiny},
  sumber/.style={draw, rounded corners, fill=green!10, text width=2.45cm, minimum height=0.9cm, align=center, font=\tiny},
  panah/.style={-{Stealth[length=2mm]}, thick}
]
  \node[tanya] (t1) at (0,0) {API apa yang ada, dan boleh dipakai siapa?};
  \node[tanya] (t2) at (2.85,0) {Bagaimana mencapai panggilan pertama hari ini?};
  \node[tanya] (t3) at (5.7,0) {Apa arti tiap jawaban dan galat?};
  \node[tanya] (t4) at (8.55,0) {Ada gangguan? Apa yang akan berubah?};

  \node[kelompok] (k1) at (0,-0.95) {Temukan};
  \node[kelompok] (k2) at (2.85,-0.95) {Mulai};
  \node[kelompok] (k3) at (5.7,-0.95) {Bangun};
  \node[kelompok] (k4) at (8.55,-0.95) {Operasikan};

  \node[halaman] (h1) at (0,-2.2) {Katalog: pemilik, audiens, versi, status\\[2pt] Syarat tiap tingkat akses};
  \node[halaman] (h2) at (2.85,-2.2) {Pendaftaran\\[2pt] Kredensial uji\\[2pt] Panduan mulai cepat};
  \node[halaman] (h3) at (5.7,-2.2) {Referensi operasi\\[2pt] Panduan galat\\[2pt] Koleksi contoh permintaan};
  \node[halaman] (h4) at (8.55,-2.2) {Status layanan\\[2pt] Jadwal versi dan \textit{sunset}\\[2pt] Dukungan dan tiket};

  \foreach \tanya/\grup/\isi in {t1/k1/h1, t2/k2/h2, t3/k3/h3, t4/k4/h4} {
    \draw[panah] (\tanya) -- (\grup);
    \draw[panah] (\grup) -- (\isi);
  }

  \node[sumber] (s1) at (0,-3.95) {Katalog dan tahap siklus hidup Bab 10};
  \node[sumber] (s2) at (2.85,-3.95) {Tangga keterbukaan Bab 3, 20 pasien sintetis};
  \node[sumber] (s3) at (5.7,-3.95) {Kontrak OpenAPI dan model galat Bab 9};
  \node[sumber] (s4) at (8.55,-3.95) {Pemantauan dan kebijakan transisi Bab 10};

  \draw[panah] (s1) -- (h1);
  \draw[panah] (s2) -- (h2);
  \draw[panah] (s3) -- (h3);
  \draw[panah] (s4) -- (h4);

  \node[font=\tiny, align=center, text width=11cm] at (4.27,-4.72)
    {Tiap isi punya satu sumber dan satu pemilik, sehingga papan informasi dan petugas layanan memberi jawaban yang sama};
\end{tikzpicture}
